%% kapa: exam=1 round=19 year=2008 period=2 form=A subject=economics subject_ko=경제원론 seq=25 number=25
%% source: Q-Net 감정평가사 기출문제 (cst00309) artlSeq=1001962 file=extracted/1008/1001962/2008년 제19회 감정평가사 1차시험 기출문제/2교시.hwp
%% conversion: kordoc 4.24.51 -> tools/gen1.py (auto); review=proofread-sonnet; flags=none
\begin{question}{25}
두 재화 $x,~y$를 소비하는 소비자의 효용함수는 $u(x,y)=x^{ 1/2 } y^{ 1/2 }$이다. 이 소비자가 효용을 극대화한다고 할 때 다음 중 옳은 것은?
\begin{choices}
  \item 소득이 2배가 되면 각 재화의 소비량이 2배보다 작게 증가한다.
  \item 소득이 2배가 되면 효용이 2배보다 작게 증가한다.
  \item 소득소비곡선은 수평선이다.
  \item 가격소비곡선은 원점을 지나는 45\textsuperscript{0}선이다.
  \item $x$재와 $y$재의 가격이 같다면 $x$재와 $y$재의 소비량은 동일하다.
\end{choices}
\end{question}
\answer{5}
